% Part: incompleteness
% Chapter: arithmetization-syntax
% Section: proofs-in-lk

\documentclass[../../../include/open-logic-section]{subfiles}

\begin{document}

\olfileid{inc}{art}{plk}
\olsection{\usetoken{P}{derivation} في $\Log{LK}$}

\begin{explain}
لكي نؤرْثِم !!{derivation}s، لا بد أن نمثل !!{derivation}s بأعداد. وبما أن
!!{derivation}s أشجار من التتابعات، يحمل كل استدلال فيها أيضًا وسمًا،
فإن التمثيل العودي هو النهج الأوضح: نمثل !!{derivation} بصفّ، مكوّناته
التتابع الختامي والوسم وتمثيلات الـ!!{derivation}s الجزئية المؤدية إلى
مقدمات الاستدلال الأخير.
\end{explain}

\begin{defn}
إذا كانت $\Gamma$ متتالية منتهية من !!{sentence}s، حيث
$\Gamma=\tuple{!A_1, \dots, !A_n}$، فإن
$\Gn{\Gamma}=\tuple{\Gn{!A_1}, \dots, \Gn{!A_n}}$.

إذا كان $\Gamma \Sequent \Delta$ تتابعًا، فإن عدد غودل لـ
$\Gamma \Sequent \Delta$ هو
\[
\Gn{\Gamma \Sequent \Delta} = \tuple{\Gn{\Gamma}, \Gn{\Delta}}
\]

إذا كان $\pi$ !!a{derivation} في $\Log{LK}$، فإن $\Gn{\pi}$ يُعرّف
كما يأتي:
\begin{enumerate}
\item إذا كان $\pi$ لا يتألف إلا من التتابع الابتدائي
  $\Gamma \Sequent \Delta$، فإن $\Gn{\pi}$ هو
  \[
    \tuple{0, \Gn{\Gamma \Sequent \Delta}}.
  \]
\item إذا انتهى $\pi$ باستدلال له مقدمة واحدة أو مقدمتان، وكانت
  $\Gamma \Sequent \Delta$ نتيجته، وكانت $\pi_1$ و$\pi_2$ البرهانين
  الجزئيين المباشرين المنتهيين بمقدمتي الاستدلال الأخير، فإن $\Gn{\pi}$ هو
  \begin{align*}
    & \tuple{1, \Gn{\pi_1}, \Gn{\Gamma \Sequent \Delta}, k} \text{ أو}\\ 
    & \tuple{2, \Gn{\pi_1}, \Gn{\pi_2}, \Gn{\Gamma \Sequent \Delta}, k}, 
  \end{align*}
  على الترتيب، حيث يعطي الجدول الآتي قيمة $k$ بحسب القاعدة المستعملة
  في الاستدلال الأخير:

  \begin{tabular}{lccccccc}
    \text{القاعدة:} & \LeftR{\Weakening} & \RightR{\Weakening} &
    \LeftR{\Contraction} & \RightR{\Contraction} &
    \LeftR{\Exchange} & \RightR{\Exchange} \\
    $k$: & 1 & 2 & 3 & 4 & 5 & 6 \\[2ex]
    \text{القاعدة:} & \LeftR{\lnot} & \RightR{\lnot} &
    \LeftR{\land} & \RightR{\land} & 
    \LeftR{\lor} & \RightR{\lor} \\
    $k$: & 7 & 8 & 9 & 10 & 11 & 12 \\[2ex]
    \text{القاعدة:} & \LeftR{\lif} & \RightR{\lif} &
    \LeftR{\lforall} & \RightR{\lforall} &
    \LeftR{\lexists} & \RightR{\lexists} \\
    $k$: & 13 & 14 & 15 & 16 & 17 & 18 \\[2ex]
    \text{القاعدة:} & \Cut & = \\
    $k$: & 19 & 20
  \end{tabular}
\end{enumerate}
\end{defn}

\begin{ex}
  تأمل !!{derivation} بسيطًا جدًا
  \begin{prooftree}
    \Axiom$!A \fCenter !A$
    \RightLabel{\LeftR{\land}}
    \UnaryInf$!A \land !B \fCenter !A$
    \RightLabel{\RightR{\lif}}
    \UnaryInf$\fCenter (!A \land !B) \lif !A$
  \end{prooftree}
  يكون عدد غودل للتتابع الابتدائي $p_0=\tuple{0,
    \Gn{!A \Sequent !A}}$. ويكون عدد غودل للـ!!{derivation} المنتهي
    بنتيجة~$\LeftR{\land}$ هو $p_1=
    \tuple{1,p_0,\Gn{!A \land !B \Sequent !A},9}$ ($1$ لأن
    $\LeftR{\land}$ لها مقدمة واحدة، ثم عدد غودل للنتيجة
    $!A \land !B \Sequent !A$، و$9$ هو العدد الذي يرمّز
    $\LeftR{\land}$). ومن ثم يكون عدد غودل للـ!!{derivation} بأكمله
    $\tuple{1,p_1,\Gn{\Sequent (!A \land !B) \lif !A},14}$، أي
  \[
  \tuple{1, \tuple{1, \tuple{0, \Gn{!A \Sequent !A}}, \Gn{!A \land !B \Sequent !A}, 9},
    \Gn{\Sequent (!A \land !B) \lif !A}, 14}.
  \]
\end{ex}

\begin{explain}
بعد تثبيت تمثيل لـ!!{derivation}s، علينا أيضًا أن نبيّن إمكان التعامل
مع !!{derivation}s كهذه عوديًا بدائيًا والتعبير بالطريقة نفسها عن خصائصها وعلاقاتها
الأساسية. بعض العمليات بسيطة؛ فمثلًا إذا أُعطينا عدد غودل~$p$
لـ!!a{derivation}، فإن $\fn{EndSeq}(p)=(p)_{(p)_0+1}$ يعطينا عدد غودل
لتتابعه الختامي، وتعطينا $\fn{LastRule}(p)=(p)_{(p)_0+2}$ شفرة
قاعدته الأخيرة. والخاصية $\fn{Sequent}(s)$ المعرّفة بـ
\[
  \len{s} = 2 \land \bforall{i<\len{(s)_0} + \len{(s)_1}}{\fn{Sent}(((s)_0 \concat (s)_1)_i)}
\]
تصدق في $s$ إذا وفقط إذا كان $s$ عدد غودل لتتابع مؤلف من
!!{sentence}s. وبعض العمليات أصعب بكثير. وسنرسم على الأقل كيفية
معالجتها. والغاية أن نثبت أن علاقة «$\pi$ !!a{derivation} لـ$!A$ من
$\Gamma$» علاقة عودية بدائية في عددي غودل لـ$\pi$ و$!A$.
\end{explain}

\begin{prop}\ollabel{prop:followsby}
  الخاصية $\fn{Correct}(p)$، التي تصدق إذا وفقط إذا كان الاستدلال
  الأخير في !!{derivation}~$\pi$ ذي عدد غودل~$p$ صحيحًا، عودية
  بدائية.
\end{prop}

\begin{proof}
  يكون $\Gamma \Sequent \Delta$ تتابعًا ابتدائيًا إذا وُجدت
  !!a{sentence}~$!A$ بحيث كان $\Gamma \Sequent \Delta$ هو
  $!A \Sequent !A$، أو وُجد حد~$t$ بحيث كان $\Gamma \Sequent \Delta$
  هو $\emptyset \Sequent \eq[t][t]$. وبلغة أعداد غودل، تصدق
  $\fn{InitSeq}(s)$ إذا وفقط إذا
  \begin{align*}
  \bexists{x < s}{} (\fn{Sent}(x) & \land
  s = \tuple{\tuple{x},\tuple{x}})
  \lor {}\\
  \bexists{t<s}{} (\fn{Term}(t) & \land
  s = \tuple{0, \tuple{\Gn{{\eq}(} \concat t \concat \Gn{,} \concat t \concat \Gn{)}}}).
  \end{align*}

  وعلينا أيضًا أن نثبت، لكل قاعدة استدلال~$R$، أن العلاقة
  $\fn{FollowsBy}_R(p)$ عودية بدائية؛ وتصدق $\fn{FollowsBy}_R(p)$
  إذا وفقط إذا كان $p$
  عدد غودل لـ!!{derivation}~$\pi$، وكان التتابع الختامي لـ$\pi$ ينتج من
  الـ!!{derivation}s الجزئية المباشرة لـ$\pi$ بتطبيق صحيح للقاعدة~$R$.

  ومن الحالات البسيطة قاعدة $\RightR{\land}$. فإذا انتهى $\pi$
  باستدلال صحيح وفق $\RightR{\land}$، كان شكلها كما يأتي:
  \begin{prooftree}
    \AxiomC{}
    \RightLabel{$\pi_1$}
    \Deduce$\Gamma \fCenter \Delta, !A$
    
    \AxiomC{}
    \RightLabel{$\pi_2$}
    \Deduce$\Gamma \fCenter \Delta, !B$

    \RightLabel{\RightR\land}
    \BinaryInf$\Gamma \fCenter \Delta, !A \land !B$
  \end{prooftree}
  وهكذا يكون الاستدلال الأخير في !!{derivation}~$\pi$ تطبيقًا صحيحًا
  لـ$\RightR{\land}$ إذا وفقط إذا وُجدت متتاليتان من !!{sentence}s،
  $\Gamma$ و$\Delta$، و!!{sentence}s عددها اثنتان، هما $!A$ و$!B$، بحيث يكون
  التتابع الختامي لـ$\pi_1$ هو $\Gamma \Sequent \Delta,!A$، والتتابع
  الختامي لـ$\pi_2$ هو $\Gamma \Sequent \Delta,!B$، والتتابع الختامي
  لـ$\pi$ هو $\Gamma \Sequent \Delta,!A \land !B$. ولا علينا إلا أن
  نترجم ذلك إلى أعداد غودل. فإذا كان
  $s=\Gn{\Gamma \Sequent \Delta}$، فإن $(s)_0=\Gn{\Gamma}$ و
  $(s)_1=\Gn{\Delta}$. ومن ثم تصدق
  $\fn{FollowsBy}_{\RightR{\land}}(p)$ إذا وفقط إذا
\begin{align*}
& \bexists{g < p}{\bexists{d < p}{\bexists{a < p}{\bexists{b < p}{\quad}}}} \\
& \qquad \fn{EndSeq}(p) = 
  \tuple{g, d \concat 
    \tuple{\Gn{(} \concat a \concat \Gn{\land} \concat b \concat \Gn{)}}} \land {} \\
& \qquad \fn{EndSeq}((p)_1) = \tuple{g, d \concat \tuple{a}} \land {}\\
& \qquad \fn{EndSeq}((p)_2) = \tuple{g, d \concat \tuple{b}} \land {}\\
& \qquad (p)_0 = 2 \land \fn{LastRule}(p) = 10.
\end{align*}
تعبّر الأسطر، على الترتيب، عن «وجود متتالية~($\Gamma$) عدد غودل لها
$g$، ومتتالية~($\Delta$) عدد غودل لها $d$، و!!a{formula}~($!A$) عدد
غودل لها $a$، و!!a{formula}~($!B$) عدد غودل لها $b$»، بحيث «يكون
التتابع الختامي لـ$\pi$ هو $\Gamma \Sequent \Delta,!A \land !B$»،
و«التتابع الختامي لـ$\pi_1$ هو $\Gamma \Sequent \Delta,!A$»،
و«التتابع الختامي لـ$\pi_2$ هو $\Gamma \Sequent \Delta,!B$»، وأن
«لـ$\pi$ اشتقاقين جزئيين مباشرين وقاعدة الاستدلال الأخيرة هي
$\RightR\land$ (ذات العدد~$10$)».

يكون الاستدلال الأخير في~$\pi$ تطبيقًا صحيحًا لـ$\RightR\lexists$
إذا وفقط إذا وُجدت متتاليتان $\Gamma$ و$\Delta$، و!!a{formula}~$!A$،
ومتغير~$x$، وحد مغلق~$t$، بحيث يكون التتابع الختامي لـ$\pi$ هو
$\Gamma \Sequent \Delta,\lexists[x][!A]$، ويكون التتابع الختامي
لـ$\pi_1$ هو $\Gamma \Sequent \Delta,\Subst{!A}{t}{x}$. ومن ثم، بلغة
أعداد غودل، تصدق $\fn{FollowsBy}_{\RightR\lexists}(p)$ إذا وفقط إذا
\begin{align*}
  & \bexists{g<p}{\bexists{d<p}{\bexists{a<p}{\bexists{x<p}{\bexists{t<p}{\quad}}}}}\\
  & \qquad \fn{Var}(x) \land \fn{ClTerm}(t) \land {}\\
  & \qquad \fn{EndSeq}(p) = 
  \tuple{
    g,
    d \concat \tuple{\Gn{\lexists} \concat x \concat a}
  } \land {}\\
  & \qquad \fn{EndSeq}((p)_1) = 
  \tuple{
    g,
    d \concat \tuple{\fn{Subst}(a, t, x)}
  } \land {}\\
  & \qquad (p)_0 = 1 \land \fn{LastRule}(p) = 18.
\end{align*}
ثم نعرّف $\fn{Correct}(p)$ بالصيغة
\begin{multline*}
  \fn{Sequent}(\fn{EndSeq}(p)) \land {}\\ 
  [(\fn{LastRule}(p) = 1 \land 
  \fn{FollowsBy}_{\LeftR\Weakening}(p)) \lor \dots \lor {}\\
  (\fn{LastRule}(p) = 20 \land \fn{FollowsBy}_{\eq}(p)) \lor {}\\
  (p)_0 = 0 \land \fn{InitSeq}(\fn{EndSeq}(p))]
\end{multline*}
يضمن السطر الأول أن التتابع الختامي لـ$p$ تتابع بالفعل مؤلف من
!!{sentence}s. ويغطي السطر الأخير الحالة التي لا يكون فيها $p$ إلا
تتابعًا ابتدائيًا.
\end{proof}

\begin{prob}
  عرّف الخصائص الآتية كما في
  \olref[inc][art][plk]{prop:followsby}:
  \begin{enumerate}
  \item $\fn{FollowsBy}_{\Cut}(p)$,
  \item $\fn{FollowsBy}_{\LeftR\lif}(p)$,
  \item $\fn{FollowsBy}_{\eq}(p)$,
  \item $\fn{FollowsBy}_{\RightR{\lforall}}(p)$.
  \end{enumerate}
  وفي الخاصية الأخيرة سيلزمك أيضًا أن تبيّن إمكان الاختبار عوديًا
  بدائيًا مما إذا كان الاستدلال الأخير في !!{derivation} ذي عدد
  غودل~$p$ يحقق شرط المتغير المميز؛ أي إن المتغير المميز~$a$ لقاعدة
  $\RightR{\lforall}$ لا يرد في التتابع الختامي.
  \end{prob}
  

\begin{prop}
  \ollabel{prop:deriv}
  العلاقة $\fn{Deriv}(p)$، التي تصدق إذا كان $p$ عدد غودل
  لـ!!{derivation} صحيح~$\pi$، عودية بدائية.
\end{prop}

\begin{proof}
  يكون !!^a{derivation}~$\pi$ صحيحًا إذا كان كل استدلال فيه تطبيقًا
  صحيحًا لقاعدة، أي إذا انتهت كل الـ!!{derivation}s الجزئية فيه باستدلال
  صحيح. ومن ثم تصدق $\fn{Deriv}(p)$ إذا وفقط إذا
  \[
  \bforall{i<\len{\fn{SubtreeSeq}(p)}}{\fn{Correct}((\fn{SubtreeSeq}(p))_i)}.
  \]
\end{proof}


\begin{prop}
افترض أن $\Gamma$ مجموعة عودية بدائية من !!{sentence}s. عندئذ تكون
العلاقة $\Prf[\Gamma](x,y)$، التي تعبّر عن أن «$x$ شفرة
!!a{derivation}~$\pi$ لـ$\Gamma_0 \Sequent !A$ من أجل مجموعة منتهية ما
$\Gamma_0\subseteq\Gamma$، وأن $y$ عدد غودل لـ$!A$»، عودية بدائية.
\end{prop}

\begin{proof}
افترض أن المحمول العودي البدائي~$R_\Gamma(y)$ يعبّر عن
«$y\in\Gamma$». علينا أن نثبت أن $\Prf[\Gamma](x,y)$، التي تصدق إذا
وفقط إذا كان $y$ عدد غودل لجملة~$!A$ وكان $x$ شفرة !!{derivation} في
$\Log{LK}$ تتابعه الختامي $\Gamma_0 \Sequent !A$، عودية بدائية.

بحسب القضية السابقة، فإن الخاصية $\fn{Deriv}(x)$، التي تصدق إذا وفقط
إذا كان $x$ شفرة !!{derivation} صحيح~$\pi$ في $\Log{LK}$، عودية
بدائية. وإذا كان $x$ شفرة كهذه، فإن $\fn{EndSeq}(x)$ شفرة التتابع
الختامي لـ$\pi$؛ ومن ثم يكون $(\fn{EndSeq}(x))_0$ شفرة طرفه الأيسر،
و$(\fn{EndSeq}(x))_1$ شفرة طرفه الأيمن. ولذلك يمكننا التعبير عن أن
«الطرف الأيمن من التتابع الختامي لـ$\pi$ هو~$!A$» بالصيغة
$\len{(\fn{EndSeq}(x))_1}=1 \land ((\fn{EndSeq}(x))_1)_0=y$. أما الطرف
الأيسر من التتابع الختامي لـ$\pi$ فمنتهٍ تلقائيًا؛ ولا يلزمنا إلا أن
نعبّر عن انتماء كل جملة فيه إلى~$\Gamma$. وهكذا نعرّف
$\Prf[\Gamma](x,y)$ بالصيغة
\begin{align*}
\Prf[\Gamma](x, y) \defiff {}&
\fn{Deriv}(x) \land {} \\
& \bforall{i <
  \len{(\fn{EndSeq}(x))_0}}{R_\Gamma(((\fn{EndSeq}(x))_0)_i)} \land {}\\
& \len{(\fn{EndSeq}(x))_1} = 1 \land ((\fn{EndSeq}(x))_1)_0 = y.
\end{align*}
\end{proof}

\end{document}
